\chapter{3次元Euclid空間}
    \section{約束}
        \begin{enumerate}
            \item $\varepsilon_{ijk}$ は Levi-Civita 記号である
        \end{enumerate}
    \section{3 次元 Euclid 空間と直交座標系の関係}
        \subsection{計量係数と単位ベクトル}
            \label{計量係数と単位ベクトル}
            $T:\realNumbers^3\to\EuclideanSpace^3$ は可逆であり $\mathrm{C}^1$ 級であるとする。
            $T$ が定める座標系を $S_T$ とする。
            $T$ のJacobi行列の列は互いに直交するものとする。
            $\bm{u}\in\realNumbers^3,\vr = T(\bm{u})$ とする。
            例えば球座標系では $\bm{u} = (r,\theta,\varphi)$ であり，$\vr = \bm{i}_1 r_1 + \bm{i}_2 r_2 + \bm{i}_3 r_3,(r_1,r_2,r_3) = (r\sin\theta\cos\varphi,r\sin\theta\sin\varphi,r\cos\theta)$ である。
            次式で定める量を座標系 $S_T$ の計量係数(Lam\'{e} coefficient)と呼ぶ。
            \[ h_i \coloneq \norm{\partDeriv{\vr}{u_i}{}}_2 \quad (i=1,2,3) \]
            $h_i$ は $T$ の Jacobi 行列 $J$ の第 $i$ 列のノルムであり，逆写像定理から $J$ は可逆なので $h_i > 0$ である。
            \par
            次式で定める単位ベクトルを $S_T$ の $u_i$ 方向単位ベクトルと呼ぶ。
            \[ \bm{j}_i := \frac{1}{h_i}\partDeriv{\vr}{u_i}{}\in\EuclideanSpace^3 \]
            $T$ のJacobi行列の列が互いに直交するため，$\bm{j}_1,\bm{j}_2,\bm{j}_3$ は互いに直交する。
            $H\coloneq\diag{h_1,h_2,h_3}$ とすると $J^\top J = H^2$ であり，$J^{-1} = H^{-2}J^\top$ である。
            つまり次式が成り立つ。
            \begin{equation}
                \partDeriv{u_i}{r_j}{} = J^{-1}[i,j] = \parens*{H^{-2}J^\top}[i,j] = \frac{1}{h_j^2}J[j,i] = \frac{1}{h_i^2}\partDeriv{r_j}{u_i}{}
                \label{直交座標系の Jacobi 行列と逆行列の要素の関係}
            \end{equation}
        \subsection{単位ベクトルの変換}
            \label{単位ベクトルの変換}
            \begin{shadebox}
                既出の記号の定義を引き継ぐ。
                3 次元 Euclid 空間の標準単位ベクトル $\bm{i}_1,\bm{i}_2,\bm{i}_3$ は直交座標系の単位ベクトル $\bm{j}_1,\bm{j}_2,\bm{j}_3$ を用いて次のように表せる。
                \[ \bm{i}_i = \sum_{j=1}^3 \frac{1}{h_j}\partDeriv{r_i}{u_j}{}\bm{j}_j \quad (i=1,2,3) \]
                つまり行列を用いて次のように表せる。
                \[ [\bm{i}_1, \bm{i}_2, \bm{i}_3] = [\bm{j}_1, \bm{j}_2, \bm{j}_3]JH^{-1} \]
                ここに $J$ は $T$ のJacobi行列であり，$H := \diag{h_1,h_2,h_3}$ である。
            \end{shadebox}
            \begin{proof}
                \quad\par
                $\bm{j}_1,\bm{j}_2,\bm{j}_3$ は正規直交基底であるから次式が成り立つ。
                \[ \bm{i}_i = \sum_{j=1}^3 (\bm{i}_i\cdot\bm{j}_j)\bm{j}_j = \sum_{j=1}^3 \parens*{\bm{i}_i\cdot\frac{1}{h_j}\partDeriv{\vr}{u_j}{}}\bm{j}_j = \sum_{j=1}^3 \parens*{\bm{i}_i\cdot\frac{1}{h_j}\sum_{k=1}^3\partDeriv{r_k}{u_j}{}\bm{i}_k}\bm{j}_j \\
                = \sum_{j=1}^3 \frac{1}{h_j}\partDeriv{r_i}{u_j}{}\bm{j}_j \]
            \end{proof}
        \subsection{ベクトルの成分の変換}
            \label{ベクトルの成分の変換}
            \begin{shadebox}
                既出の記号の定義を引き継ぐ。
                3次元Euclid空間上のベクトル$\vA(\vr) = \bm{i}_1 A_1(\vr) + \bm{i}_2 A_2(\vr) + \bm{i}_3 A_3(\vr)$を$u_1,u_2,u_3$直交座標系で表したものを$\bm{j}_1 \tilde{A}_1(\bm{u}) + \bm{j}_2 \tilde{A}_2(\bm{u}) + \bm{j}_3 \tilde{A}_3(\bm{u})$とするとき，$A_1,A_2,A_3$と$\tilde{A}_1,\tilde{A}_2,\tilde{A}_3$の間には次の関係式が成り立つ。
                \[ [A_1,A_2,A_3]^\top = JH^{-1}[\tilde{A}_1,\tilde{A}_2,\tilde{A}_3]^\top \]
                ここに$J$は$T$のJacobi行列であり，$H := \diag{h_1,h_2,h_3}$である。
            \end{shadebox}
            \begin{proof}
                \begin{align*}
                    \phantom{=} &[\bm{i}_1,\bm{i}_2,\bm{i}_3][A_1,A_2,A_3]^\top = \vA(\vr) = \bm{j}_1 \tilde{A}_1 + \bm{j}_2 \tilde{A}_2 + \bm{j}_3 \tilde{A}_3 = [\bm{j}_1,\bm{j}_2,\bm{j}_3][\tilde{A}_1,\tilde{A}_2,\tilde{A}_3]^\top \\
                    = &[\bm{i}_1,\bm{i}_2,\bm{i}_3]JH^{-1}[\tilde{A}_1,\tilde{A}_2,\tilde{A}_3]^\top
                \end{align*}
            \end{proof}
        \subsection{勾配}
            \label{勾配}
            \begin{shadebox}
                既出の記号の定義を引き継ぐ。
                $\bm{u}\in\realNumbers^3,\vr=T(\bm{u})\in\EuclideanSpace^3$ とする。
                $\phi:\EuclideanSpace^3\to\complexNumbers$は$\EuclideanSpace^3$上で$\mathrm{C}^1$ 級であるとする。
                $\tilde{\phi}(\bm{u}) \coloneq \phi(\vr)$ とする。
                $\phi$ の勾配 $\nabla_\vr\phi(\vr)$ は $\bm{u}$ を用いて次式で表せる。
                % \[ \nabla_\vr\phi(\vr) = \sum_{i=1}^3\frac{\bm{j}_i}{h_i}\partDeriv{\tilde{\phi}}{u_i}{} = \bracks*{\frac{\bm{j}_1}{h_1}, \frac{\bm{j}_2}{h_2}, \frac{\bm{j}_3}{h_3}}^\top\nabla_\bm{u}\tilde{\phi}(\bm{u}) \]
                \[ \nabla_\vr\phi(\vr) = \sum_{i=1}^3\frac{\bm{j}_i}{h_i}\partDeriv{\tilde{\phi}}{u_i}{} \]
            \end{shadebox}
            \begin{proof}
                % 線形代数を過用した証明。長くて読みにくい。
                % \begin{gather*}
                %     \bm{j}_i = \frac{1}{h_i}\partDeriv{\vr}{u_i}{} = \frac{1}{h_i}\partDerivLong{\sum_{j=1}^3 r_j\bm{i}_j}{u_i}{} = \frac{1}{h_i}[\bm{i}_1,\bm{i}_2,\bm{i}_3]\bracks*{\partDeriv{r_1}{u_i}{},\partDeriv{r_2}{u_i}{},\partDeriv{r_3}{u_i}{}}^\top \\
                %     \therefore\;[\bm{j}_1,\bm{j}_2,\bm{j}_3] = [\bm{i}_1,\bm{i}_2,\bm{i}_3]\bracks*{\frac{1}{h_j}\partDeriv{r_i}{u_j}{}}_{(i,j)\in\{1,2,3\}^2} = [\bm{i}_1,\bm{i}_2,\bm{i}_3]JH^{-1}
                % \end{gather*}
                % ここに$\quad H \coloneq \diag{h_1,h_2,h_3}$であり，$J$は$T$のJacobi行列である。
                % 以上より次式が成り立つ。
                % \[ [\bm{i}_1,\bm{i}_2,\bm{i}_3] = [\bm{j}_1,\bm{j}_2,\bm{j}_3]HJ^{-1} \]
                % $\partDeriv{\phi}{r_i}{}$を評価すると次式を得る。
                % \begin{gather*}
                %     \partDeriv{\phi}{r_i}{} = \sum_{j=1}^3\partDeriv{\tilde{\phi}}{u_j}{}\partDeriv{u_j}{r_i}{} \\
                %     \begin{split}
                %         \therefore \quad \nabla_\vr\phi(\vr) &= \bracks*{\partDeriv{\phi}{r_1}{}, \partDeriv{\phi}{r_2}{}, \partDeriv{\phi}{r_3}{}}^\top = \bracks*{\partDeriv{u_j}{r_i}{}}_{(i,j)\in\{1,2,3\}^2}\bracks*{\partDeriv{\tilde{\phi}}{u_1}{},\partDeriv{\tilde{\phi}}{u_2}{},\partDeriv{\tilde{\phi}}{u_3}{}}^\top \\
                %         &= \parens*{J^{-1}}^\top\nabla_\bm{u}\tilde{\phi}(\bm{u})
                %     \end{split}
                % \end{gather*}
                % 最後の等号は$\bm{u} = T^{-1}(\vr)$と逆関数定理による。
                % 以上より，次式が成り立つ。
                % \[ \nabla_\vr\phi(\vr) = [\bm{j}_1,\bm{j}_2,\bm{j}_3]HJ^{-1}\parens*{J^{-1}}^\top\nabla_\bm{u}\tilde{\phi}(\bm{u}) \]
                % $J^{-1}\parens*{J^{-1}}^\top = J^{-1}\parens*{J^\top}^{-1} = \parens*{J^\top J}^{-1}$であり，$T$に関する仮定から$J$の列は互いに直交し，計量係数の定義から$J^\top J$の第$i$対角成分は$h_i^2$である。
                % よって$\parens*{J^\top J}^{-1} = H^{-2}$であり，次式が成り立つ。
                % \[ \nabla_\vr\phi(\vr) = [\bm{j}_1,\bm{j}_2,\bm{j}_3]H^{-1}\nabla_\bm{u}\tilde{\phi}(\bm{u}) = \bracks*{\frac{\bm{j}_1}{h_1}, \frac{\bm{j}_2}{h_2}, \frac{\bm{j}_3}{h_3}}^\top\nabla_\bm{u}\tilde{\phi}(\bm{u}) \]
                % (2026-05-04) 洗練
                \quad\par
                \ref{単位ベクトルの変換} と合成関数の微分の連鎖律から次式が成り立つ。
                \begin{align*}
                    \nabla_\vr\phi(\vr) &= \sum_{i=1}^3 \bm{i}_i\partDeriv{\tilde{\phi}(\bm{u}(\vr))}{r_i}{} = \sum_{i=1}^3\parens*{\sum_{j=1}^3 \bm{j}_j\frac{1}{h_j}\partDeriv{r_i}{u_j}{}}\parens*{\sum_{k=1}^3 \partDeriv{\tilde{\phi}}{u_k}{}\partDeriv{u_k}{r_i}{}} \\
                    &= \sum_{j=1}^3 \bm{j}_j\frac{1}{h_j}\sum_{k=1}^3 \partDeriv{\tilde{\phi}}{u_k}{}\sum_{i=1}^3 \partDeriv{u_k}{r_i}{}\partDeriv{r_i}{u_j}{} = \sum_{j=1}^3 \bm{j}_j\frac{1}{h_j}\sum_{k=1}^3 \partDeriv{\tilde{\phi}}{u_k}{}\parens*{J^{-1}J}_{k,j} \\
                    &= \sum_{j=1}^3 \bm{j}_j\frac{1}{h_j}\sum_{k=1}^3 \partDeriv{\tilde{\phi}}{u_k}{}\delta_{k,j} = \sum_{j=1}^3 \bm{j}_j\frac{1}{h_j}\partDeriv{\tilde{\phi}}{u_j}{}
                \end{align*}
            \end{proof}
    \subsection{発散}
        \label{発散}
        \begin{shadebox}
            既出の記号の定義を引き継ぐ。
            $\bm{u}\in\realNumbers^3,\vr=T(\bm{u})\in\EuclideanSpace^3$ とする。
            さらに $T$ は $\mathrm{C}^2$ 級であるとし，$\partDeriv{h_i}{u_i}{} = 0\;(i=1,2,3)$ が成り立つとする(例: 円柱座標系, 球座標系)。
            3 次元 Euclid 空間上の $\mathrm{C}^1$ 級のベクトル場 $\vA(\vr) = \bm{i}_1 A_1(\vr) + \bm{i}_2 A_2(\vr) + \bm{i}_3 A_3(\vr)$ を $u_1,u_2,u_3$ 直交座標系で表したものを $\bm{j}_1 \tilde{A}_1(\bm{u}) + \bm{j}_2 \tilde{A}_2(\bm{u}) + \bm{j}_3 \tilde{A}_3(\bm{u})$ とするとき，次式が成り立つ。
            \[ \nabla_\vr\cdot\vA(\vr) = \frac{1}{h_1 h_2 h_3}\bracks*{ \partDeriv{h_2 h_3 \tilde{A}_1(\bm{u})}{u_1}{} + \partDeriv{h_3 h_1 \tilde{A}_2(\bm{u})}{u_2}{} + \partDeriv{h_1 h_2 \tilde{A}_3(\bm{u})}{u_3}{} } = \frac{1}{h_1 h_2 h_3}\sum_{i=1}^3 \partDerivLong{\frac{h_1 h_2 h_3 \tilde{A}_i(\bm{u})}{h_i}}{u_i}{} \]
        \end{shadebox}
        \begin{proof}
            \begin{align*}
                \nabla_\vr\cdot\vA(\vr) &= \sum_{i=1}^3 \partDeriv{A_i(\vr)}{r_i}{} = \sum_{i=1}^3 \sum_{j=1}^3 \partDeriv{A_i(\vr)}{u_j}{}\partDeriv{u_j}{r_i}{} = \sum_{i=1}^3 \sum_{j=1}^3 \parens*{\partDerivLong{\sum_{k=1}^3\frac{1}{h_k}\partDeriv{r_i}{u_k}{}\tilde{A}_k(\bm{u})}{u_j}{}}\partDeriv{u_j}{r_i}{} \\
                &= \sum_{i=1}^3 \sum_{j=1}^3 \parens*{\partDerivLong{\sum_{k=1}^3\frac{1}{h_k}\partDeriv{r_i}{u_k}{}\tilde{A}_k(\bm{u})}{u_j}{}}\frac{1}{h_j^2}\partDeriv{r_i}{u_j}{}\quad(\because\ref{直交座標系の Jacobi 行列と逆行列の要素の関係}) \tag{1}
            \end{align*}
            \par
            式(1)の$\tilde{A}_k(\bm{u})$に関する項に着目すると，以下のように計算される。
            但し以下の式で$\delta_{i,j}$はKroneckerのデルタを表す。
            \begin{align*}
                \phantom{=} &\sum_{i=1}^3 \sum_{j=1}^3 \parens*{\partDerivLong{\frac{1}{h_k}\partDeriv{r_i}{u_k}{}\tilde{A}_k(\bm{u})}{u_j}{}}\frac{1}{h_j^2}\partDeriv{r_i}{u_j}{} = \sum_{j=1}^3 \frac{1}{h_j^2} \sum_{i=1}^3 \parens*{\partDerivLong{\frac{1}{h_k}\partDeriv{r_i}{u_k}{}\tilde{A}_k(\bm{u})}{u_j}{}}\partDeriv{r_i}{u_j}{} \\
                = &\sum_{j=1}^3 \frac{1}{h_j^2} \sum_{i=1}^3 \bracks*{ \parens*{\partDerivLong{\frac{1}{h_k}}{u_j}{}}\tilde{A}_k(\bm{u})\partDeriv{r_i}{u_k}{}\partDeriv{r_i}{u_j}{} + \frac{1}{h_k}\partDeriv{\tilde{A}_k(\bm{u})}{u_j}{}\partDeriv{r_i}{u_k}{}\partDeriv{r_i}{u_j}{} + \frac{\tilde{A}_k(\bm{u})}{h_k}\partDerivIIHetero{r_i}{u_j}{u_k}\partDeriv{r_i}{u_j}{} } \\
                = &\sum_{j=1}^3 \frac{1}{h_j^2} \bracks*{ \parens*{\partDerivLong{\frac{1}{h_k}}{u_j}{}}\tilde{A}_k(\bm{u})\partDeriv{\vr}{u_k}{}\cdot\partDeriv{\vr}{u_j}{} + \frac{1}{h_k}\partDeriv{\tilde{A}_k(\bm{u})}{u_j}{}\partDeriv{\vr}{u_k}{}\cdot\partDeriv{\vr}{u_j}{} + \frac{\tilde{A}_k(\bm{u})}{h_k}\partDerivIIHetero{\vr}{u_j}{u_k}\cdot\partDeriv{\vr}{u_j}{} } \\
                = &\sum_{j=1}^3 \frac{1}{h_j^2} \bracks*{ \parens*{\partDerivLong{\frac{1}{h_k}}{u_j}{}}\tilde{A}_k(\bm{u})h_k^2\delta_{k,j} + \frac{1}{h_k}\partDeriv{\tilde{A}_k(\bm{u})}{u_j}{}h_k^2\delta_{k,j} + \frac{\tilde{A}_k(\bm{u})}{h_k}\partDerivIIHetero{\vr}{u_j}{u_k}\cdot\partDeriv{\vr}{u_j}{} } \\
                = &\frac{1}{h_k^2} \bracks*{ \parens*{\partDerivLong{\frac{1}{h_k}}{u_k}{}}\tilde{A}_k(\bm{u})h_k^2 + h_k\partDeriv{\tilde{A}_k(\bm{u})}{u_k}{} } + \frac{\tilde{A}_k(\bm{u})}{h_k} \sum_{j=1}^3 \frac{1}{h_j^2} \partDerivIIHetero{\vr}{u_j}{u_k}\cdot\partDeriv{\vr}{u_j}{} \\
                = &\frac{1}{h_k}\partDeriv{\tilde{A}_k(\bm{u})}{u_k}{} + \frac{\tilde{A}_k(\bm{u})}{h_k} \sum_{j=1}^3 \underbrace{\frac{1}{h_j^2}\partDerivIIHetero{\vr}{u_j}{u_k}\cdot\partDeriv{\vr}{u_j}{}}_{(2)} \quad \parens*{\because \text{仮定より}\partDerivLong{\frac{1}{h_k}}{u_j}{} = 0} \tag{3}
            \end{align*}
            式(2)を評価する。
            $j=k$のとき次式より0である。
            \[ \partDeriv{\vr}{u_k}{2}\cdot\partDeriv{\vr}{u_k}{} = \frac{1}{2}\partDerivLong{\parens*{\partDeriv{\vr}{u_k}{}\cdot\partDeriv{\vr}{u_k}{}}}{u_k}{} = \frac{1}{2}\partDeriv{h_k^2}{u_k}{} = 0 \]
            $j\neq k$のとき次式を得る。
            \begin{align*}
                \phantom{=} &\frac{1}{h_j^2}\partDerivIIHetero{\vr}{u_j}{u_k}\cdot\partDeriv{\vr}{u_j}{} = \frac{1}{h_j^2}\partDerivIIHetero{\vr}{u_k}{u_j}\cdot\partDeriv{\vr}{u_j}{} \\
                &= \frac{1}{h_j^2}\times\frac{1}{2} \partDerivLong{\parens*{\partDeriv{\vr}{u_j}{}\cdot\partDeriv{\vr}{u_j}{}}}{u_k}{} = \frac{1}{h_j^2}\times\frac{1}{2}\partDeriv{h_j^2}{u_k}{} = \frac{1}{h_j}\partDeriv{h_j}{u_k}{}
            \end{align*}
            これらを式(3)に適用して次式を得る。
            % 但し以下で$[i+j]$は$i+j$を3で割った余りを表す。
            \begin{align*}
                (3) &= \frac{1}{h_k}\partDeriv{\tilde{A}_k(\bm{u})}{u_k}{} + \frac{\tilde{A}_k(\bm{u})}{h_k}\bracks*{ \frac{1}{h_{(k-1)\%3+1}}\partDeriv{h_{(k-1)\%3+1}}{u_k}{} + \frac{1}{h_{(k-1)\%3+2}}\partDeriv{h_{(k-1)\%3+2}}{u_k}{} } \\
                &= \frac{1}{h_1 h_2 h_3}\partDerivLong{\frac{h_1 h_2 h_3\tilde{A}_k(\bm{u})}{h_k}}{u_k}{}
            \end{align*}
        \end{proof}
    % \subsection{回転} (2026-05-04) 証明断念
    %     \begin{shadebox}
    %         既出の記号の定義を引き継ぐ。
    %         $\bm{u}\in\realNumbers^3,\vr=T(\bm{u})\in\EuclideanSpace^3$ とする。
    %         さらに$T$は$\mathrm{C}^2$級であるとし，次式が成り立つとする(例: 円柱座標系, 球座標系)。
    %         \[ \partDerivIIHetero{\vr}{u_j}{u_i}\cdot\partDeriv{\vr}{u_k}{} = 0 \quad \parens*{\{i,j,k\} = \{1,2,3\}} \tag{1} \]
    %         3次元Euclid空間上の$\mathrm{C}^1$級のベクトル場$\vA(\vr) = \bm{i}_1 A_1(\vr) + \bm{i}_2 A_2(\vr) + \bm{i}_3 A_3(\vr)$を$u_1,u_2,u_3$直交座標系で表したものを$\bm{j}_1 \tilde{A}_1(\bm{u}) + \bm{j}_2 \tilde{A}_2(\bm{u}) + \bm{j}_3 \tilde{A}_3(\bm{u})$とするとき，次式が成り立つ。
    %         \begin{align*}
    %             \nabla_\vr\times\vA(\vr) &= \frac{\bm{j}_1}{h_2 h_3}\bracks*{ \partDeriv{h_3\tilde{A}_3}{u_2}{} - \partDeriv{h_2\tilde{A}_2}{u_3}{} } + \frac{\bm{j}_2}{h_3 h_1}\bracks*{ \partDeriv{h_1\tilde{A}_1}{u_3}{} - \partDeriv{h_3\tilde{A}_3}{u_1}{} } + \frac{\bm{j}_3}{h_1 h_2}\bracks*{ \partDeriv{h_2\tilde{A}_2}{u_1}{} - \partDeriv{h_1\tilde{A}_1}{u_2}{} } \\
    %             &= \sum_{i=1}^3 \bm{j}_i \varepsilon_{ijk}\frac{1}{h_j h_k}\partDeriv{h_k \tilde{A}_k}{u_j}{}
    %         \end{align*}
    %     \end{shadebox}
    %     \begin{proof}
    %         \begin{align*}
    %             \phantom{=} &\nabla_\vr\times\vA(\vr) = \nabla_\vr \times \sum_{i=1}^3 A_i(\vr)\bm{i}_i = \sum_{i=1}^3 \nabla_\vr \times A_i(\vr)\bm{i}_i = \sum_{i=1}^3 (\nabla_\vr A_i(\vr))\times\bm{i}_i\quad (\because \text{\ref{ベクトルのスカラー関数倍の回転}}) \\
    %             = &\sum_{i=1}^3 \sum_{j=1}^3 \frac{\bm{j}_j}{h_j}\partDeriv{A_i(\vr)}{u_j}{}\times\bm{i}_i = \sum_{i=1}^3 \sum_{j=1}^3 \frac{\bm{j}_j}{h_j}\parens*{\partDerivLong{\sum_{k=1}^3 \frac{\tilde{A}_k(\bm{u})}{h_k}\partDeriv{r_i}{u_k}{}}{u_j}{}} \times \sum_{l=1}^3 \frac{\bm{j}_l}{h_l}\partDeriv{r_i}{u_l}{} \\
    %             &(\because \text{\ref{勾配}, \ref{ベクトルの成分の変換}, \ref{単位ベクトルの変換}})
    %         \end{align*}
    %         以下では表記の簡便さのため，$\tilde{A}_k(\bm{u})$を単に$\tilde{A}_k$と書く。
    %         上式の第$m\in\{1,2,3\}$成分を評価する。
    %         第$m$成分に寄与するのは$(j,l)=\parens*{[m+1],[m+2]},\parens*{[m+2],[m+1]}$のときであり，計算すると次式を得る。
    %         \begin{align*}
    %             \phantom{=} &\sum_{i=1}^3 \left[ \frac{1}{h_{[m+1]}}\parens*{\partDerivLong{\sum_{k=1}^3 \frac{\tilde{A}_k(\bm{u})}{h_k}\partDeriv{r_i}{u_k}{}}{u_{[m+1]}}{}}\frac{1}{h_{[m+2]}}\partDeriv{r_i}{u_{[m+2]}}{} \right. \\
    %             &\left. - \frac{1}{h_{[m+2]}}\parens*{\partDerivLong{\sum_{k=1}^3 \frac{\tilde{A}_k(\bm{u})}{h_k}\partDeriv{r_i}{u_k}{}}{u_{[m+2]}}{}}\frac{1}{h_{[m+1]}}\partDeriv{r_i}{u_{[m+1]}}{} \right] \\
    %             &= \frac{1}{h_{[m+1]}h_{[m+2]}}\sum_{k=1}^3\bracks*{ \sum_{i=1}^3\parens*{\partDerivLong{\frac{\tilde{A}_k}{h_k}\partDeriv{r_i}{u_k}{}}{u_{[m+1]}}{}}\partDeriv{r_i}{u_{[m+2]}}{} - \sum_{i=1}^3\parens*{\partDerivLong{\frac{\tilde{A}_k}{h_k}\partDeriv{r_i}{u_k}{}}{u_{[m+2]}}{}}\partDeriv{r_i}{u_{[m+1]}}{} } \tag{2}
    %         \end{align*}
    %         この式の[ ]内第1項を評価すると次式を得る。
    %         \begin{align*}
    %             \phantom{=} &\sum_{k=1}^3\sum_{i=1}^3\parens*{\partDerivLong{\frac{\tilde{A}_k}{h_k}\partDeriv{r_i}{u_k}{}}{u_{[m+1]}}{}}\partDeriv{r_i}{u_{[m+2]}}{} \\
    %             = &\sum_{k=1}^3 \bracks*{ \parens*{\partDerivLong{\frac{\tilde{A}_k}{h_k}}{u_{[m+1]}}{}}\sum_{i=1}^3\partDeriv{r_i}{u_k}{}\partDeriv{r_i}{u_{[m+2]}}{} + \frac{\tilde{A}_k}{h_k}\sum_{i=1}^3\partDerivIIHetero{r_i}{u_{[m+1]}}{u_k}\partDeriv{r_i}{u_{[m+2]}}{} } \\
    %             = &\sum_{k=1}^3 \bracks*{ \parens*{\partDerivLong{\frac{\tilde{A}_k}{h_k}}{u_{[m+1]}}{}}\partDeriv{\vr}{u_k}{}\cdot\partDeriv{\vr}{u_{[m+2]}}{} + \frac{\tilde{A}_k}{h_k}\partDerivIIHetero{\vr}{u_{[m+1]}}{u_k}\cdot\partDeriv{\vr}{u_{[m+2]}}{} } \\
    %             = &\parens*{\partDerivLong{\frac{\tilde{A}_{[m+2]}}{h_{[m+2]}}}{u_{[m+1]}}{}}h_{[m+2]}^2 + \sum_{k\in\{[m+1],[m+2]\}} \frac{\tilde{A}_k}{h_k}\partDerivIIHetero{\vr}{u_{[m+1]}}{u_k}\cdot\partDeriv{\vr}{u_{[m+2]}}{}
    %         \end{align*}
    %         最後の等号成立には式(1)を用いた。
    %         式(2)の[ ]内第2項も同様に評価すると，式(2)の$1/(h_{[m+1]}h_{[m+2]})$の右側の部分は次式になる。
    %         \begin{align*}
    %             \phantom{=} &\parens*{\partDerivLong{\frac{\tilde{A}_{[m+2]}}{h_{[m+2]}}}{u_{[m+1]}}{}}h_{[m+2]}^2 - \parens*{\partDerivLong{\frac{\tilde{A}_{[m+1]}}{h_{[m+1]}}}{u_{[m+2]}}{}}h_{[m+1]}^2 \\
    %             \phantom{=}    & + \sum_{k\in\{[m+1],[m+2]\}}\frac{\tilde{A}_k}{h_k}\parens*{\partDerivIIHetero{\vr}{u_{[m+1]}}{u_k}\cdot\partDeriv{\vr}{u_{[m+2]}}{} - \partDerivIIHetero{\vr}{u_{[m+2]}}{u_k}\cdot\partDeriv{\vr}{u_{[m+1]}}{}} \tag{3}
    %         \end{align*}
    %         この式の第3項の()内を評価する。
    %         $T$は$\mathrm{C}^2$級であると仮定しているから，微分順序の入れ替えが可能であることに留意する。
    %         $k=[m+1]$のときは次式になる。
    %         \begin{align*}
    %             \phantom{=} &\partDeriv{\vr}{u_{[m+1]}}{2}\cdot\partDeriv{\vr}{u_{[m+2]}}{} - \partDerivIIHetero{\vr}{u_{[m+2]}}{u_{[m+1]}}\cdot\partDeriv{\vr}{u_{[m+1]}}{} \\
    %             = &\partDerivLong{\parens*{\partDeriv{\vr}{u_{[m+1]}}{}\cdot\partDeriv{\vr}{u_{[m+2]}}{}}}{u_{[m+1]}}{} - \partDeriv{\vr}{u_{[m+1]}}{}\cdot\partDerivIIHetero{\vr}{u_{[m+1]}}{u_{[m+2]}}{} - \partDerivIIHetero{\vr}{u_{[m+2]}}{u_{[m+1]}}\cdot\partDeriv{\vr}{u_{[m+1]}}{} \\
    %             = &-\partDeriv{\vr}{u_{[m+1]}}{}\cdot\partDerivIIHetero{\vr}{u_{[m+2]}}{u_{[m+1]}}{} - \partDerivIIHetero{\vr}{u_{[m+2]}}{u_{[m+1]}}\cdot\partDeriv{\vr}{u_{[m+1]}}{} \\
    %             = &-\partDerivLong{\parens*{\partDeriv{\vr}{u_{[m+1]}}{}\cdot\partDeriv{\vr}{u_{[m+1]}}{}}}{u_{[m+2]}}{} = -\partDeriv{h_{[m+1]}^2}{u_{[m+2]}}{}
    %         \end{align*}
    %         $k=[m+2]$のときは次式になる。
    %         \begin{align*}
    %             \phantom{=} &\partDerivIIHetero{\vr}{u_{[m+1]}}{u_{[m+2]}}\cdot\partDeriv{\vr}{u_{[m+2]}}{} - \partDeriv{\vr}{u_{[m+2]}}{2}\cdot\partDeriv{\vr}{u_{[m+1]}}{} \\
    %             = &\partDerivIIHetero{\vr}{u_{[m+1]}}{u_{[m+2]}}\cdot\partDeriv{\vr}{u_{[m+2]}}{} - \partDerivLong{\parens*{\partDeriv{\vr}{u_{[m+2]}}{}\cdot\partDeriv{\vr}{u_{[m+1]}}{}}}{u_{[m+2]}}{} + \partDeriv{\vr}{u_{[m+2]}}{}\cdot\partDerivIIHetero{\vr}{u_{[m+2]}}{u_{[m+1]}} \\
    %             = &\partDerivIIHetero{\vr}{u_{[m+1]}}{u_{[m+2]}}\cdot\partDeriv{\vr}{u_{[m+2]}}{} + \partDeriv{\vr}{u_{[m+2]}}{}\cdot\partDerivIIHetero{\vr}{u_{[m+1]}}{u_{[m+2]}} \\
    %             = &\partDerivLong{\parens*{\partDeriv{\vr}{u_{[m+2]}}{}\cdot\partDeriv{\vr}{u_{[m+2]}}{}}}{u_{[m+1]}}{} = \partDeriv{h_{[m+2]}^2}{u_{[m+1]}}{}
    %         \end{align*}
    %         これらの式を(3)に適用すると次式を得る。
    %         \begin{align*}
    %             (3) &= \parens*{\partDerivLong{\frac{\tilde{A}_{[m+2]}}{h_{[m+2]}}}{u_{[m+1]}}{}}h_{[m+2]}^2 + \frac{\tilde{A}_{[m+2]}}{h_{[m+2]}}\partDeriv{h_{[m+2]}^2}{u_{[m+1]}}{} - \parens*{\partDerivLong{\frac{\tilde{A}_{[m+1]}}{h_{[m+1]}}}{u_{[m+2]}}{}}h_{[m+1]}^2 - \frac{\tilde{A}_{[m+1]}}{h_{[m+1]}}\partDeriv{h_{[m+1]}^2}{u_{[m+2]}}{} \\
    %             &= \partDerivLong{\frac{\tilde{A}_{[m+2]}}{h_{[m+2]}}h_{[m+2]}^2}{u_{[m+1]}}{} - \partDerivLong{\frac{\tilde{A}_{[m+1]}}{h_{[m+1]}}h_{[m+1]}^2}{u_{[m+2]}}{} \\
    %             &= \partDeriv{h_{[m+2]}\tilde{A}_{[m+2]}}{u_{[m+1]}}{} - \partDeriv{h_{[m+1]}\tilde{A}_{[m+1]}}{u_{[m+2]}}{}
    %         \end{align*}
    %         これを式(2)に適用して定理の主張を得る。
    %     \end{proof}
    \section{Laplacian}
        \begin{shadebox}
            既出の記号の定義を引き継ぐ。
            $\bm{u}\in\realNumbers^3,\vr=T(\bm{u})\in\EuclideanSpace^3$ とする。
            さらに$T$は$\mathrm{C}^2$級であるとし，$\partDeriv{h_i}{u_i}{} = 0\;(i=1,2,3)$が成り立つとする(例: 円柱座標系, 球座標系)。
            $\phi:\realNumbers^3\to\complexNumbers$は$\realNumbers^3$上で$\mathrm{C}^2$級であるとする。
            $\tilde{\phi}(\bm{u}) \coloneq \phi(T(\bm{u}))$とする。
            $\nabla^2_\vr\phi(\vr)$は$\bm{u}$を用いて次式で表せる。
            \[ \nabla^2_\vr\phi(\vr) = \frac{1}{h_1 h_2 h_3}\bracks*{ \partDerivLong{\parens*{\frac{h_2 h_3}{h_1}\partDeriv{\tilde{\phi}(\bm{u})}{u_1}{}}}{u_1}{} + \partDerivLong{\parens*{\frac{h_3 h_1}{h_2}\partDeriv{\tilde{\phi}(\bm{u})}{u_2}{}}}{u_2}{} + \partDerivLong{\parens*{\frac{h_1 h_2}{h_3}\partDeriv{\tilde{\phi}(\bm{u})}{u_3}{}}}{u_3}{} } \]
        \end{shadebox}
        \begin{proof}
            \ref{勾配},\ref{発散}を用いて容易に示せる。
        \end{proof}
    \section{諸公式}
        \subsection{ベクトル三重積の公式: \texorpdfstring{$A\times(B\times C) = \inProd{A}{C}B - \inProd{A}{B}C$}{A×(B×C) = (A・C)B - (A・B)C}}
            \begin{proof}
                \quad\par
                直交座標系をとり，単位ベクトルを $\bm{j}_1,\bm{j}_2,\bm{j}_3$ とする。
                $A,B,C$を次のようにおく。
                \begin{align*}
                    A &= \bm{j}_1 A_1 + \bm{j}_2 A_2 + \bm{j}_3 A_3 \\
                    B &= \bm{j}_1 B_1 + \bm{j}_2 B_2 + \bm{j}_3 B_3 \\
                    C &= \bm{j}_1 C_1 + \bm{j}_2 C_2 + \bm{j}_3 C_3
                \end{align*}
                すると
                \begin{align*}
                    A\times(B\times C) &= (\bm{j}_1 A_1 + \bm{j}_2 A_2 + \bm{j}_3 A_3) \times \bracks*{\bm{j}_1(B_2 C_3 - B_3 C_2) + \bm{j}_2(B_3 C_1 - B_1 C_3) + \bm{j}_3(B_1 C_2 - B_2 C_1)} \\
                    &= \phantom{+}\bm{j}_3 A_1(B_3 C_1 - B_1 C_3) - \bm{j}_2 A_1(B_1 C_2 - B_2 C_1) \\
                    &\phantom{=} -\bm{j}_3 A_2(B_2 C_3 - B_3 C_2) + \bm{j}_1 A_2(B_1 C_2 - B_2 C_1) \\
                    &\phantom{=} +\bm{j}_2 A_3(B_2 C_3 - B_3 C_2) - \bm{j}_1 A_3(B_3 C_1 - B_1 C_3) \\
                    &= \phantom{+}\bm{j}_1 \bracks*{(A_2 C_2 + A_3 C_3)B_1 - (A_2 B_2 + A_3 B_3)C_1} \\
                    &\phantom{=} + \bm{j}_2 \bracks*{(A_1 C_1 + A_3 C_3)B_2 - (A_1 B_1 + A_3 B_3)C_2} \\
                    &\phantom{=} + \bm{j}_3 \bracks*{(A_1 C_1 + A_2 C_2)B_3 - (A_1 B_1 + A_2 B_2)C_3} \\
                    &= \phantom{+}\bm{j}_1 \bracks*{(\inProd{A}{C} - A_1 C_1)B_1 - (\inProd{A}{B} - A_1 B_1)C_1} \\
                    &\phantom{=} + \bm{j}_2 \bracks*{(\inProd{A}{C} - A_2 C_2)B_2 - (\inProd{A}{B} - A_2 B_2)C_2} \\
                    &\phantom{=} + \bm{j}_3 \bracks*{(\inProd{A}{C} - A_3 C_3)B_3 - (\inProd{A}{B} - A_3 B_3)C_3} \\
                    &= \phantom{+}\bm{j}_1 \bracks*{\inProd{A}{C}B_1 - \inProd{A}{B}C_1} \\
                    &\phantom{=} + \bm{j}_2 \bracks*{\inProd{A}{C}B_2 - \inProd{A}{B}C_2} \\
                    &\phantom{=} + \bm{j}_3 \bracks*{\inProd{A}{C}B_3 - \inProd{A}{B}C_3} \\
                    &= \inProd{A}{C}B - \inProd{A}{B}C
                \end{align*}
            \end{proof}
        \subsection{\texorpdfstring{$\inProd{\bm{a}\times\bm{b}}{\bm{c}\times\bm{d}} = \inProd{\bm{a}}{\bm{c}}\inProd{\bm{b}}{\bm{d}} - \inProd{\bm{a}}{\bm{d}}\inProd{\bm{b}}{\bm{c}}$}{(a×b)・(c×d) = (a・c)(b・d) - (a・d)(b・c)}}
            \begin{proof}
                \quad\par
                ベクトル三重積の公式を利用する。
                \[ \inProd{\bm{a}\times\bm{b}}{\bm{c}\times\bm{d}} = \inProd{\bm{a}}{\bm{b}\times(\bm{c}\times\bm{d})} = \inProd{a}{\inProd{\bm{d}}{\bm{b}}\bm{c} - \inProd{\bm{b}}{\bm{c}}\bm{d}} = \inProd{\bm{a}}{\bm{c}}\inProd{\bm{b}}{\bm{d}} - \inProd{\bm{a}}{\bm{d}}\inProd{\bm{b}}{\bm{c}} \]
            \end{proof}
        \subsection{(系) \texorpdfstring{$\inProd{\bm{a}\times\bm{b}}{\bm{a}\times\bm{c}} = \norm{\bm{a}}^2\inProd{\bm{b}}{\bm{c}} - \inProd{\bm{a}}{\bm{b}}\inProd{\bm{a}}{\bm{c}}$}{(a×b)・(a×c) = ||a||\string^2(b・c) - (a・b)(a・c)}}
            %実は直前の公式よりもこちらの方を先に思いついた。直前の公式から直ちに導かれるが，当時思いついた証明を書き残しておく。
            \begin{proof}
                直前の公式から直ちに導ける。
                % \begin{align*}
                %     \inProd{\bm{a}\times\bm{b}}{\bm{a}\times\bm{c}} &= (a_2b_1 - a_1b_2) (a_2c_1 - a_1c_2) + (a_3b_1 - a_1b_3) (a_3c_1 - a_1c_3) + (a_3b_2 - a_2b_3) (a_3c_2 - a_2c_3) \\
                %     &=  {a_1}^2 b_2 c_2 + {a_1}^2 b_3 c_3 - a_1 a_2 b_1 c_2 - a_1 a_2 b_2 c_1 - a_1 a_3 b_1 c_3 - a_1 a_3 b_3 c_1 + {a_2}^2 b_1 c_1 \\
                %     &\quad + {a_2}^2 b_3 c_3 - a_2 a_3 b_2 c_3 - a_2 a_3 b_3 c_2 + {a_3}^2 b_1 c_1 + {a_3}^2 b_2 c_2 \\
                %     &= ({a_1}^2 + {a_2}^2 + {a_3}^2)(b_1c_1 + b_2c_2 + b_3c_3) - {a_1}^2b_1c_1 - {a_2}^2b_2c_2 - {a_3}^2b_3c_3 \\
                %     &\quad - a_1 a_2 b_1 c_2 - a_1 a_2 b_2 c_1 - a_1 a_3 b_1 c_3 - a_1 a_3 b_3 c_1 - a_2 a_3 b_2 c_3 - a_2 a_3 b_3 c_2 \\
                %     &= ({a_1}^2 + {a_2}^2 + {a_3}^2)(b_1c_1 + b_2c_2 + b_3c_3) - (a_1b_1 + a_2b_2 + a_3b_3)(a_1c_1 + a_2c_2 + a_3c_3) \\
                %     &= \norm{\bm{a}}^2\inProd{\bm{b}}{\bm{c}} - \inProd{\bm{a}}{\bm{b}}\inProd{\bm{a}}{\bm{c}}
                % \end{align*}
            \end{proof}
        \subsection{発散の別表現: \texorpdfstring{$\nabla\cdot A = \frac{1}{|V|}\lim_{|V|\to 0} \integrate{\partial V}{}{A(\vr)\cdot\bm{n}(\vr)}{2}{r}$}{発散はベクトル場の法線面積分の体積の商の極限である}}
            \begin{shadebox}
                $\vr_0\in\EuclideanSpace^3$とし，$\Omega\subset\EuclideanSpace^3$上のベクトル場$A$が$\vr_0$の近傍で$C^1$級であるとする。
                $\vr_0$を含む微小な領域を$V$とし，$\partial V$を$V$の境界，$|V|$を$V$の体積とする。
                但し$\partial V$は$C^1$級であるとする。
                このとき次式が成り立つ。
                \[ (\nabla\cdot A)(\vr_0) = \frac{1}{|V|}\lim_{|V|\to 0}\integrate{\partial V}{}{A(\vr)\cdot\bm{n}(\vr)}{2}{r} \]
                ここに$\bm{n}(\vr)$は$\vr\in\partial V$に於ける外向きの法線ベクトルである。
            \end{shadebox}
            \begin{proof}
                \quad\par
                発散定理より次式が成り立つ。
                \[ \integrate{\partial V}{}{A(\vr)\cdot\bm{n}(\vr)}{2}{r} = \integrate{V}{}{(\nabla\cdot A)(\vr)}{3}{r} \]
                $\varepsilon>0$を任意に小さくとる。
                $A$が$\vr_0$の近傍で$C^1$級であるという仮定から，$V$を充分小さくとれば，$|(\nabla\cdot A)(\vr) - (\nabla\cdot A)(\vr_0)| < \varepsilon$である。
                よって次式が成り立つ。
                \[ \abs{ \frac{1}{|V|}\integrate{V}{}{(\nabla\cdot A)(\vr)}{3}{r} - (\nabla\cdot A)(\vr_0) } < \varepsilon \]
                $|V|\to 0$とすれば定理の主張が成り立つ。
            \end{proof}
        \subsection{ベクトルのスカラー関数倍の発散\texorpdfstring{: $\nabla_\vr\cdot f(\vr)\vA(\vr) = \parens*{\nabla_\vr f(\vr)}\cdot \vA(\vr) + f(\vr)\nabla_\vr\cdot\vA(\vr)$}{}}
            \begin{shadebox}
                $\vr\in\EuclideanSpace^3$とする。
                $f:\EuclideanSpace^3\to\complexNumbers,A:\EuclideanSpace^3\to\EuclideanSpace^3$はともに$\mathrm{C}^1$級とする。
                座標系の取り方に依らず次式が成り立つ。
                \[ \nabla_\vr\cdot f(\vr)\vA(\vr) = \parens*{\nabla_\vr f(\vr)}\cdot \vA(\vr) + f(\vr)\nabla_\vr\cdot\vA(\vr) \]
            \end{shadebox}
            \begin{proof}
                \quad\par
                \ref{空間と座標系の関係}の\ref{Euclid 空間上の，座標が陽に現れない関係式が成り立つ必要十分条件} より，直交座標系上で成立を示せば十分である。
                $\bm{u}\in\realNumbers^3,\vr=T(\bm{u}),\tilde{f}(\bm{u})\coloneq f(\vr),\tilde{A}(\bm{u})\coloneq A(\vr)$ とする。
                \begin{align*}
                    \nabla_\vr\cdot f(\vr)\vA(\vr) &= \frac{1}{h_1 h_2 h_3}\sum_{i=1}^3 \partDerivLong{\frac{h_1 h_2 h_3}{h_i}\tilde{f} \tilde{A}_i}{u_i}{} = \frac{1}{h_1 h_2 h_3}\sum_{i=1}^3\parens*{\partDeriv{\tilde{f}}{u_i}{}\frac{h_1 h_2 h_3}{h_i} \tilde{A}_i + \tilde{f}\partDerivLong{\frac{h_1 h_2 h_3}{h_i} \tilde{A}_i}{u_i}{}} \\
                    &= \sum_{i=1}^3 \frac{1}{h_i}\partDeriv{\tilde{f}}{u_i}{}\tilde{A}_i + \tilde{f}\frac{1}{h_1 h_2 h_3}\sum_{i=1}^3 \partDerivLong{\frac{h_1 h_2 h_3}{h_i}\tilde{A}_i}{u_i}{} = (\nabla f)\cdot\vA + f\nabla\cdot\vA
                \end{align*}
            \end{proof}
        \subsection{ベクトルのスカラー関数倍の回転\texorpdfstring{: $\nabla_\vr\times f(\vr)\vA(\vr) = \parens*{\nabla_\vr f(\vr)}\times \vA(\vr) + f(\vr)\nabla_\vr\times\vA(\vr)$}{}}
            \label{ベクトルのスカラー関数倍の回転}
            \begin{shadebox}
                $\vr\in\EuclideanSpace^3$とする。
                $f:\EuclideanSpace^3\to\complexNumbers,A:\EuclideanSpace^3\to\EuclideanSpace^3$はともに$\mathrm{C}^1$級とする。
                座標系の取り方に依らず次式が成り立つ。
                \[ \nabla_\vr\times f(\vr)\vA(\vr) = \parens*{\nabla_\vr f(\vr)}\times \vA(\vr) + f(\vr)\nabla_\vr\times\vA(\vr) \]
            \end{shadebox}
            \begin{proof}
                \quad\par
                \ref{空間と座標系の関係}の\ref{Euclid 空間上の，座標が陽に現れない関係式が成り立つ必要十分条件} より，直交座標系上で成立を示せば十分である。
                $\bm{u}\in\realNumbers^3,\vr=T(\bm{u}),\tilde{f}(\bm{u})\coloneq f(\vr),\tilde{A}(\bm{u})\coloneq A(\vr)$ とする。
                \begin{align*}
                    \nabla_\vr\times f(\vr)\vA(\vr) &= \sum_{i=1}^3 \bm{j}_i\varepsilon_{ijk}\frac{1}{h_j h_k}\partDerivLong{h_k \tilde{f} \tilde{A}_k}{u_j}{} = \sum_{i=1}^3 \bm{j}_i\varepsilon_{ijk}\frac{1}{h_j h_k}\parens*{\partDeriv{\tilde{f}}{u_j}{}h_k \tilde{A}_k + \tilde{f}\partDeriv{h_k \tilde{A}_k}{u_j}{}} \\
                    &= \sum_{i=1}^3 \bm{j}_i\varepsilon_{ijk}\frac{1}{h_j}\partDeriv{\tilde{f}}{u_j}{} \tilde{A}_k + \tilde{f}\sum_{i=1}^3 \bm{j}_i\varepsilon_{ijk}\partDeriv{h_k \tilde{A}_k}{u_j}{} = (\nabla f)\times\vA + f\nabla\times\vA
                \end{align*}
            \end{proof}
        \subsection{距離のn乗の勾配\texorpdfstring{: $\nabla_\vr\norm{\vr+\bm{a}}^n = n\norm{\vr+\bm{a}}^{n-2}(\vr+\bm{a})$}{}}
            \begin{shadebox}
                $\bm{a}\in\EuclideanSpace^3,\vr\in\EuclideanSpace^3$とする。
                座標系の取り方に依らず次式が成り立つ。
                \[ \nabla_\vr\norm{\vr+\bm{a}}^n = n\norm{\vr+\bm{a}}^{n-2}(\vr+\bm{a}) \]
            \end{shadebox}
            \begin{proof}
                \quad\par
                \ref{空間と座標系の関係}の\ref{Euclid 空間上の，座標が陽に現れない関係式が成り立つ必要十分条件} より，直交座標系上で成立を示せば十分である。
                $\bm{u}\in\realNumbers^3,\vr=T(\bm{u})$ とする。
                \begin{align*}
                    \nabla_\vr\norm{\vr+\bm{a}}^n &= \sum_{i=1}^3 \bm{j}_i \frac{1}{h_i}\partDerivLong{\parens*{\norm{\vr+\bm{a}}^2}^{n/2}}{u_i}{} = \sum_{i=1}^3 \bm{j}_i \frac{1}{h_i}\frac{n}{2}\parens*{\norm{\vr+\bm{a}}^2}^{n/2-1} 2(r_i+a_i)\partDeriv{r_i}{u_i}{} \\
                    &= n\norm{\vr+\bm{a}}^{n-2} \sum_{i=1}^3 \bm{j}_i \frac{1}{h_i}(r_i+a_i)h_i = n\norm{\vr+\bm{a}}^{n-2}(\vr+\bm{a})
                \end{align*}
            \end{proof}
    \section{Vector Laplacian}
        \label{VectorLaplacian}
        \begin{definition}
            $\vA:\vr\in\EuclideanSpace^3 \mapsto \vA(\vr) \in \EuclideanSpace^3$($\vA$の各成分は$\mathrm{C}^2$級)の Vector Laplacian $\nabla^2 \vA$は次式で定義される。
            \[ \nabla^2 \vA \coloneq \nabla(\nabla\cdot \vA) - \nabla\times(\nabla\times \vA) \]
        \end{definition}
    \section{無発散ベクトル場のベクトル・ポテンシャルの存在}
        \begin{shadebox}
            $\vA:\EuclideanSpace^3\to\EuclideanSpace^3$ の発散が 0 であるとき，適当な $\bm{P}:\EuclideanSpace^3\to\EuclideanSpace^3$ が存在して $\nabla\times\bm{P} = \vA$ となる。
        \end{shadebox}
        \begin{proof}
            \quad\par
            \ref{空間と座標系の関係}の\ref{Euclid 空間上の，座標が陽に現れない関係式が成り立つ必要十分条件} より，直角座標系上で成立を示せば十分である。
            次式で構成された $\bm{P}$ が $\nabla\times\bm{P} = \vA$ を満たすことを示す。
            \begin{align*}
                P_1(\vr) &= \frac{1}{3}\parens*{\integrate{0}{r_3}{A_2(r_1,r_2,\tilde{r_3})}{}{\tilde{r_3}} - \integrate{0}{r_2}{A_3(r_1,\tilde{r_2},r_3)}{}{\tilde{r_2}}} \\
                P_2(\vr) &= \frac{1}{3}\parens*{\integrate{0}{r_1}{A_3(\tilde{r_1},r_2,r_3)}{}{\tilde{r_1}} - \integrate{0}{r_3}{A_1(r_1,r_2,\tilde{r_3})}{}{\tilde{r_3}}} \\
                P_3(\vr) &= \frac{1}{3}\parens*{\integrate{0}{r_2}{A_1(r_1,\tilde{r_2},r_3)}{}{\tilde{r_2}} - \integrate{0}{r_1}{A_2(\tilde{r_1},r_2,r_3)}{}{\tilde{r_1}}}
            \end{align*}
            $\nabla\times\bm{P}$ の第 1 成分は次式である。
            \begin{align*}
                (\nabla\times\bm{P})_1(\vr) &= \frac{1}{3}\parens*{A_1(\vr) - \integrate{0}{r_1}{\partDeriv{A_2}{r_2}{}(\tilde{r_1},r_2,r_3)}{}{\tilde{r_1}}} - \frac{1}{3}\parens*{\integrate{0}{r_1}{\partDeriv{A_3}{r_3}{}(\tilde{r_1},r_2,r_3)}{}{\tilde{r_1}} - A_1(\vr)} \\
                &= \frac{2}{3}A_1(\vr) - \frac{1}{3}\integrate{0}{r_1}{\parens*{\partDeriv{A_2}{r_2}{}(\tilde{r_1},r_2,r_3) + \partDeriv{A_3}{r_3}{}(\tilde{r_1},r_2,r_3)}}{}{\tilde{r_1}} \\
                &= \frac{2}{3}A_1(\vr) - \frac{1}{3}\integrate{0}{r_1}{-\partDeriv{A_1}{r_1}{}(\tilde{r_1},r_2,r_3)}{}{\tilde{r_1}}\quad(\because\nabla\cdot\vA = 0) \\
                &= A_1(\vr)
            \end{align*}
            同様にして他の成分も示せる。
        \end{proof}
    \section{Helmholtzの定理}
        \subsection{試み1：不完全}
            証明が不完全であるが，一定の成果を得たので書き残しておく。
            \subsubsection{補題: 全空間に渡る積分が存在するための十分条件}
                \label{全空間に渡る積分が存在するための十分条件}
                \begin{shadebox}
                    $\vr\in\realNumbers^3$とする。
                    $f:\realNumbers^3\to\complexNumbers$は十分遠方に於いて，その絶対値がある$\alpha>0$に対して$o(\norm{\vr}_2^{-3-\alpha})$であるとする。
                    このとき$\realNumbers^3$上での$f$の積分は存在する。
                    すなわち次式が成り立つ。
                    \[ R>0, V(R) \coloneq \setComprehension{\vr\in\realNumbers^3}{\norm{\vr}_2 \leq R} ,\lim_{R\to\infty} \abs{\integrate{V(R)}{}{f(\vr)}{3}{r}} < \infty \]
                \end{shadebox}
                \begin{proof}
                    \quad\par
                    $R>0$は十分大きいものとし，$I(R)$を次式で定義する。
                    \[ I(R) \coloneq \integrate{V(R)}{}{f(\vr)}{3}{r} \]
                    球座標変換$\vr = r[\sin\theta\cos\phi,\sin\theta\sin\phi,\cos\theta]^\top$を用いて次式を得る。
                    \[ I(R) = \integrate{0}{R}{\integrate{0}{\pi}{\integrate{0}{2\pi}{f(\vr)r^2\sin\theta}{}{\phi}}{}{\theta}}{}{r} \]
                    $0<R_0<R$なる$R_0$をとり，$I(R) = I(R_0) + I_2(R)$と分解する。
                    ここに$I_2(R) \coloneq I(R) - I(R_0)$である。
                    \begin{align*}
                        |I_2(R)| &\leq \integrate{R_0}{R}{\integrate{0}{\pi}{\integrate{0}{2\pi}{|f(\vr)|r^2\sin\theta}{}{\phi}}{}{\theta}}{}{r} = 4\pi\integrate{R_0}{R}{o(\norm{\vr}_2^{-3-\alpha})r^2}{}{r} \\
                        &= 4\pi\integrate{R_0}{R}{o(r^{-1-\alpha})}{}{r} < \infty \text{ as } R\to\infty
                    \end{align*}
                \end{proof}
            \subsubsection{補題: 与えられた発散を有し，回転が\texorpdfstring{$\bm{0}$}{0}であるベクトル場の構成}
                \label{与えられた発散を有し，回転が無いベクトル場の構成}
                Diracのデルタ関数を使わずに証明したいが，未達成。今後の課題。
            \subsubsection{補題: 与えられた回転を有し，発散が0であるベクトル場の構成}
                \label{与えられた回転を有し，発散が無いベクトル場の構成}
                Diracのデルタ関数を使わずに証明したいが，未達成。今後の課題。
            \subsubsection{Helmholtzの定理: 任意の\texorpdfstring{$\mathrm{C}^1$}{C\string^1}級のベクトル場であって，\ref{全空間に渡る積分が存在するための十分条件}の意味で十分遠方で0に収束するものは回転が\texorpdfstring{$\bm{0}$}{0}である\texorpdfstring{$\mathrm{C}^1$}{C\string^1}級の場と発散が0である\texorpdfstring{$\mathrm{C}^1$}{C\string^1}級の場に分解できる}
                任意に与えられた，発散および回転が十分遠方で0に収束するような$\mathrm{C}^1$級のベクトル場$A$が与えられたとき，補題\ref{与えられた発散を有し，回転が無いベクトル場の構成}により$A$と同じ発散を持ち，かつ回転が$\bm{0}$であるようなベクトル場$B$を構成できる。
                このとき$A - B$は発散が0であるので，補題\ref{与えられた回転を有し，発散が無いベクトル場の構成}により，回転がこれに一致するようなベクトル場$C$を構成できる。
                これらを用いて$A = B + C$と分解すればよい。
        \subsection{試み2：完成度高}
            これはFourier変換を用いる発想であり，\cite{qm3_chap5_j} 11.2 節の内容を筆者が理解してまとめたものである。
            \subsubsection{Helmholtzの定理: 任意の\texorpdfstring{$\mathrm{C}^1$}{C\string^1}級のベクトル場は回転が\texorpdfstring{$\bm{0}$}{0}である\texorpdfstring{$\mathrm{C}^1$}{C\string^1}級の場と発散が0である\texorpdfstring{$\mathrm{C}^1$}{C\string^1}級の場に分解できる}
            \newcommand*{\vBO}{\vB_\mathrm{O}(\vk)}
            \newcommand*{\vBP}{\vB_\mathrm{P}(\vk)}
            \begin{proof}
                \quad\par
                \ref{空間と座標系の関係}の\ref{Euclid 空間上の，座標が陽に現れない関係式が成り立つ必要十分条件} より，直角座標系上で成立を示せば十分である。
                ベクトル場 $\vA:\realNumbers^3\to\complexNumbers^3$ は $\mathrm{C}^1$ 級であるとする。
                $\vA$ のFourier変換を $\vB:\realNumbers^3\to\complexNumbers^3$ とする（次式）。
                \begin{gather*}
                    \vB(\vk) = \frac{1}{\sqrt{(2\pi)^3}}\integrate{\vr\in\realNumbers^3}{}{\vA(\vr)\exp\parens*{-i\vk\cdot\vr}}{3}{\vr} \\
                    \vA(\vr) = \frac{1}{\sqrt{(2\pi)^3}}\integrate{\vk\in\realNumbers^3}{}{\vB(\vk)\exp\parens*{i\vk\cdot\vr}}{3}{\vk}
                \end{gather*}
                これを用いて $\vA$ の発散と回転を計算すると次式を得る。
                \begin{align*}
                    \parens*{\nabla\cdot\vA}(\vr) &= \frac{1}{\sqrt{(2\pi)^3}}\integrate{\vk\in\realNumbers^3}{}{\nabla\cdot\parens*{\vB(\vk)\exp\parens*{i\vk\cdot\vr}}}{3}{\vk} \\
                    &= \frac{1}{\sqrt{(2\pi)^3}}\integrate{\vk\in\realNumbers^3}{}{\parens*{\nabla\exp\parens*{i\vk\cdot\vr}}\cdot\vB(\vk)}{3}{\vk} \\
                    &= \frac{1}{\sqrt{(2\pi)^3}}\integrate{\vk\in\realNumbers^3}{}{i\vk\cdot\vB(\vk)\exp\parens*{i\vk\cdot\vr}}{3}{\vk} \tag{1} \\
                    \parens*{\nabla\times\vA}(\vr) &= \frac{1}{\sqrt{(2\pi)^3}}\integrate{\vk\in\realNumbers^3}{}{\nabla\times\parens*{\vB(\vk)\exp\parens*{i\vk\cdot\vr}}}{3}{\vk} \\
                    &= \frac{1}{\sqrt{(2\pi)^3}}\integrate{\vk\in\realNumbers^3}{}{\parens*{\nabla\exp\parens*{i\vk\cdot\vr}}\times\vB(\vk)}{3}{\vk} \\
                    &= \frac{1}{\sqrt{(2\pi)^3}}\integrate{\vk\in\realNumbers^3}{}{i\vk\times\vB(\vk)\exp\parens*{i\vk\cdot\vr}}{3}{\vk} \tag{2}
                \end{align*}
                上式中の $\vB(\vk)$ を $\vk$ 方向の成分 $\vBP$ とそれに垂直な成分 $\vBO \coloneq \vB(\vk) - \vBP$ に分解する(`P': `parallel', `O': `orthogonal')。
                $\vBP$ を具体的に書くと次式である。
                \[ \vBP = \parens*{\frac{\vk}{\norm{\vk}_2}\cdot\vB(\vk)}\frac{\vk}{\norm{\vk}_2} = \parens*{\frac{\vk}{\norm{\vk}_2^2}\cdot\vB(\vk)}\vk \]
                これを用いて $\vA$ を分解する。
                \begin{align*}
                    \vA(\vr) &= \underbrace{\frac{1}{\sqrt{(2\pi)^3}}\integrate{\vk\in\realNumbers^3}{}{\vBP\exp\parens*{i\vk\cdot\vr}}{3}{\vk}}_{\mathrm{(A)}} + \underbrace{\frac{1}{\sqrt{(2\pi)^3}}\integrate{\vk\in\realNumbers^3}{}{\vBO\exp\parens*{i\vk\cdot\vr}}{3}{\vk}}_{\mathrm{(B)}} \tag{3}\\
                    &= \vARL(\vr) + \vADL(\vr) \quad \text{where} \quad \vARL(\vr) \coloneq \text{(A)}, \quad \vADL(\vr) \coloneq \text{(B)}
                \end{align*}
                `RL', `DL' はそれぞれ `rotation-less', `divergence-less' の意味であるが，その理由は直後に明らかになる。
                \par
                式(1),(2)の計算と全く同様にして次式を得る。
                \begin{align*}
                    \nabla\cdot\vADL(\vr) &= \frac{1}{\sqrt{(2\pi)^3}}\integrate{\vk\in\realNumbers^3}{}{i\vk\cdot\vBO\exp\parens*{i\vk\cdot\vr}}{3}{\vk} = 0 \\
                    \nabla\times\vARL(\vr) &= \frac{1}{\sqrt{(2\pi)^3}}\integrate{\vk\in\realNumbers^3}{}{i\vk\times\vBP\exp\parens*{i\vk\cdot\vr}}{3}{\vk} = \bm{0}
                \end{align*}
                これと式(3)より $(\nabla\cdot\vARL)(\vk) = \parens*{\nabla\cdot\vA}(\vr),\quad (\nabla\times\vADL)(\vk) = \parens*{\nabla\times\vA}(\vr)$ も成り立つ。
            \end{proof}
    \section{諸定理}
        \subsection{球殻定理}
            \begin{shadebox}
                $R>0,V \coloneq \setComprehension{\vr\in\EuclideanSpace^3}{\norm{\vr}_2\leq R}$とする。
                ベクトル場$\vA(\vr)$を次式で定義する。
                \[ \vA(\vr) \coloneq \integrate{V}{}{\frac{\vr-\vr'}{\norm{\vr-\vr'}_2^3}}{3}{r'} \]
                このとき次が成り立つ。
                \[
                    \vA(\vr) = \frac{4}{3}\pi
                    \begin{cases}
                        \vr & (\norm{\vr}_2\leq R) \\
                        \frac{R^3}{\norm{\vr}_2^3}\vr & (\norm{\vr}_2 > R)
                    \end{cases}
                \]
            \end{shadebox}
            \begin{proof}
                \quad\par
                必要なら系全体を回転させ，一般性を失わず$\vr = [0,0,r]^\top\;(r>0)$とする。
                球座標変換を用いて積分を実行する。
                $\vr' \coloneq r'[\sin\theta\cos\phi,\sin\theta\sin\phi,\cos\theta]^\top$とすると
                \begin{align*}
                    \vA(\vr) &= \integrate{0}{R}{\integrate{0}{\pi}{\integrate{0}{2\pi}{\frac{\vr-\vr'}{\norm{\vr-\vr'}_2^3}r'^2(\sin\theta)}{}{\phi}}{}{\theta}}{}{r'} \\
                    &= 2\pi\bm{i}_3\integrate{0}{R}{\integrate{0}{\pi}{\frac{r-r'\cos\theta}{(r^2+r'^2-2rr'\cos\theta)^{3/2}}r'^2(\sin\theta)}{}{\theta}}{}{r'}
                \end{align*}
                ここに$\bm{i}_3$は第3軸方向の単位ベクトルである。
                \par
                $f(r) \coloneq -(r^2+r'^2-2rr'\cos\theta)^{-1/2}r'^2\sin\theta$とすると次式が成り立つ。
                \[ \deriv{f}{r}{}(r) = \frac{(r-r'\cos\theta)r'^2\sin\theta}{(r^2+r'^2-2rr'\cos\theta)^{3/2}} \]
                上式右辺の$(r',\theta)=(0,0)$近傍での振る舞いを考える。
                $r'=r+\Delta r\;(\Delta r\neq 0)$とすると，$\sin\theta = O(\theta),\cos\theta = 1+O(\theta^2)$なので分母・分子ともに$r^3 O(\theta^3)$となる。
                よって積分領域に於いて$f$は$r$の$\mathrm{C}^1$級関数であるから積分と微分の順序交換が可能であり，次式が成り立つ。
                \begin{align*}
                    \vA(\vr) &= 2\pi\bm{i}_3\derivLong{\integrate{0}{R}{\integrate{0}{\pi}{f(r)}{}{\theta}}{}{r'}}{r}{} \\
                    &= -2\pi\bm{i}_3\derivLong{\integrate{0}{R}{\integrate{0}{\pi}{\frac{r'^2\sin\theta}{\sqrt{r^2+r'^2-2rr'\cos\theta}}}{}{\theta}}{}{r'}}{r}{}
                \end{align*}
                $g(\theta) \coloneq \frac{r'}{r}\sqrt{r^2+r'^2-2rr'\cos\theta}$とすると次式が成り立つ。
                \[ \deriv{g}{\theta}{}(\theta) = \frac{r'^2\sin\theta}{\sqrt{r^2+r'^2-2rr'\cos\theta}} \]
                よって次式が成り立つ。
                \[ \vA(\vr) = -2\pi\bm{i}_3\derivLong{\integrate{0}{R}{\bracks*{g(\theta)}_0^\pi}{}{r'}}{r}{} = -2\pi\bm{i}_3\derivLong{\integrate{0}{R}{\frac{r'}{r}\parens*{|r+r'|-|r-r'|}}{}{r'}}{r}{} \]
                $R<r$のときは次式が成り立つ。
                \[ \vA(\vr) = -2\pi\bm{i}_3\derivLong{\integrate{0}{R}{\frac{2r'^2}{r}}{}{r'}}{r}{} = \frac{4\pi R^3}{3r^2}\bm{i}_3 = \frac{4\pi R^3}{3}\frac{\vr}{\norm{\vr}_2^3} \]
                $R\geq r$のときは$r'$に関する積分領域を$[0,r]$と$[r,R]$に分けて考える。
                前者については$R<r$のときの結果で$R$を$r$で置き換えたものとなる。
                後者については$|r+r'|-|r-r'| = 2r$となり，次式が成り立つ。
                \[ \vA(\vr) = -2\pi\bm{i}_3\derivLong{\integrate{0}{R}{2r'}{}{r'}}{r}{} = \bm{0} \]
            \end{proof}
            これを検証したMathematicaノートブックを\texttt{parts/vectorCalculus/mathematica/\\ shell\_theorem.nb}に置いてある。
        \subsection{非負領域の共通部分の非負領域}
            \begin{shadebox}
                $\realNumbers^n$上で考える。$\vv_1,\dotsc,\vv_m\;(m\leq n)$は一次独立であるとし，$\vv$の非負領域を
                \[P(\vv)\coloneq\braces*{\bm{w}|\inProd{\vv}{\bm{w}}\geq0})\]
                で定義する。あるベクトル$\bm{u}$が$\bigcap_{i=1}^{m}{P(\vv_i)}$の非負領域にある，すなわち
                \[\inProd{\bm{u}}{\bm{w}}\geq0,\forall\bm{w}\in\bigcap_{i=1}^{m}{P(\vv_i)}\]
                であるとき，$\bm{u}$は適当な定数$a_i\geq0(i=1,\dotsc,m)$によって
                \[\bm{u} = \sum_{i=1}^{m}{a_i\vv_i}\]
                と表せる。
            \end{shadebox}
            \begin{proof}
                \quad\par
                    まず$\bm{u}$が$\vv_1,\dotsc,\vv_m$の一次結合で表せることを示す。\\
                    $W \coloneq \Span{\vv_1,\dotsc,\vv_m}$とすると$\realNumbers^n = W\oplus W^\bot$であるから，適当な定数$c_1,\dotsc,c_m$及び適当なベクトル$\bm{w}_t\in W^\bot$を用いて
                    \begin{equation}
                        \bm{u} = \sum_{i=1}^{m}{c_i\vv_i} + \bm{w}_t
                    \end{equation}
                    と表せる。$\bm{p}$を$\bigcap_{i=1}^{m}{P(\vv_i)}$の任意のベクトルとすると
                    \[\inProd{\vv_i}{\bm{p}}\geq0,\forall\vv_i\in W\]
                    また，$\bm{w}_{t2}$を$W^\bot$の任意のベクトルとすると
                    \[\inProd{\vv_i}{\bm{w}_{t2}}=0,\forall\vv_i\in W\]
                    従って$\bm{p}' \coloneq \bm{p}+\bm{w}_{t2}$とすると
                    \[\inProd{\vv_i}{\bm{p}'}\geq0,\forall\vv_i\in W\]
                    であるから$\bm{p'}$もまた$\bigcap_{i=1}^{m}{P(\vv_i)}$のベクトルである。従って$\inProd{\bm{u}}{\bm{p}'}\geq0$であるから
                    \[0\leq\inProd{\bm{u}}{\bm{p}'} = \inProd{\sum_{i=1}^{m}{c_i\vv_i}}{\bm{p}} + \inProd{\bm{w}_{t1}}{\bm{w}_{t2}}\]
                    である。然るに$\bm{w}_{t2}\in W^\bot$は任意であるから，$\bm{w}_{t1}=\bm{0}$でない限り$\inProd{\bm{w}_{t1}}{\bm{w}_{t2}}$の値をいくらでも小さくできる。
                    よって上式が成り立つためには$\bm{w}_{t1}=\bm{0}$でなくてはならならず，式(1)より結局
                    \[\bm{u} = \sum_{i=1}^{m}{c_i\vv_i}\]
                    と表せることになる。
                    \par
                    次に$c_i$が全て非負であることを示す。%そのために$\parens*{\bm{d}_i\in\bigcap_{i=1}^{m}{P(\vv_i)}}\land\inProd{\vv_j}{\bm{d}_i}=\delta_{ij}\;(\delta_{ij}\text{はKroneckerのデルタ})$なるベクトル$\bm{d}_i$が存在することを示そう。
                    $\vv_1,\dotsc,\vv_m\;(m\leq n)$は一次独立であったから次の行列
                    \[
                        A = \bracks*{
                            \begin{array}{c}
                                {\vv_1}^\mathrm{T}\\
                                {\vv_2}^\mathrm{T}\\
                                \vdots\\
                                {\vv_m}^\mathrm{T}
                            \end{array}
                        }
                    \]
                    の階数は$m$である。従って次の方程式
                    \[
                        A\vx = \bracks*{
                            \begin{array}{c}
                                0\\
                                \vdots\\
                                1\\
                                \vdots\\
                                0
                            \end{array}
                        }
                        \begin{array}{c}
                            \quad\\
                            \quad\\
                            \leftarrow i\\
                            \quad\\
                            \quad
                        \end{array}
                    \]
                    は解を持つ。この解を$\bm{d}_i$とすれば$\inProd{\vv_j}{\bm{d}_i}=\delta_{ij}$であるから$\bm{d}_i\in\bigcap_{i=1}^{m}{P(\vv_i)}$。
                    よって
                    \[0\leq\inProd{\bm{u}}{\bm{d}_i} = c_i\]
            \end{proof}
    \subsection{Stokesの定理}
        \begin{shadebox}
            $c$を空間上の任意の$C^1$級の閉曲線，$S$を，$c$を周$\partial S$とする任意の$C^2$級の曲面, $A$を$C^1$級のベクトル場とするとき
            \[
                \oint_c{A\cdot\dvc} = \iint_S{\nabla\times A}\cdot\dS
            \]
        \end{shadebox}
        \begin{proof}
            \quad\par
            $S$が$S(x,y)$の如く$x,y$の関数で表せる場合について証明すれば十分である。
            なぜなら，$S$が複雑な場合は$S(s,y),S(y,z),S(z,x)$いずれかの型の2変数関数で表せる曲面に分割して，各部分で，これから証明する$S(x,y)$型の場合と同じ手続きを踏んで最後に合計すれば全体の結果になるからである。
            領域分割を行う場合は分割境界でも線積分を行うことになるが，その回数は2で，互いに逆向きであるから，合計する段階で相殺して0になるため最終的な結果には現れない。
            \begin{figure}[H]
                \centering
                \includegraphics[keepaspectratio, scale=0.4]
                {\currfiledir/figs/Stokes-theorem/divide.pdf}
                \caption{$S$の分割}
            \end{figure}
            以降$S$は$S(x,y)$型の曲面とする。$\partial S$を$t$をパラメータとする閉曲線とする。
            すなわち
            \[\partial S\;:\;x=\phi(t), y=\psi(t), t\;:\;0\rightarrow1\]
            また，$xy$平面への$S$の射影を$\Omega$とする。
            \begin{figure}[H]
                \centering
                \includegraphics[keepaspectratio, scale=0.4]
                {\currfiledir/figs/Stokes-theorem/illust.pdf}
                \caption{状況}
            \end{figure}
            では証明しよう。
            線積分で表された定理の左辺を変形して面積分に置き換えて，定理の右辺と見比べる流れで行く。
            \par
            関数の引数を全部書いていてはスペースがいくらあっても足りないので省略して書くことになる(例えば$A_x(x,y,z)$を単に$A_x$と書いたり)のだが，こうすると独立変数と従属変数の関係を見誤る危険が生じるため注意が必要だ。
            迂闊に計算すると罠にはまる。
            特に危ない所では敢えて引数を明示する。
            \begin{align*}
                \oint_c{A\cdot\dvc} &= \oint_{\partial S}{\parens*{\vix A_x+\viy A_y+\viz A_z}\cdot\parens*{\vix \dx+\viy \dy+\viz \dz}} \\
                &= \oint_{\partial S}{\parens*{A_x\dx+A_y\dy}} + \oint_{\partial S}{A_z\dz} \\
                &= \underbrace{\oint_{\partial \textcolor{red}{\Omega}}{\parens*{A_x(x,y,S(x,y))\dx+A_y(x,y,S(x,y))\dy}}}_{\ctext{A}} + \underbrace{\oint_{\partial S}{A_z\dz}}_{\ctext{B}}
            \end{align*}
            \ctext{A}で積分経路が$\partial S$から$\partial\Omega$に変更されたことに注意。Greenの定理を使いたいから$xy$平面に落としたのだ。
            Greenの定理より
            \[
                \ctext{A} = \iint_\Omega{\bracks*{\underbrace{\partDerivLong{A_y(x,y,S(x,y))}{x}{}}_{\ctext{C}}-\underbrace{\partDerivLong{A_x(x,y,S(x,y))}{y}{}}_{\ctext{D}}}\dx\dy}
            \]
            \ctext{C}が罠だ。引数を明示していないと$\left.\partDeriv{A_y}{x}{}\right|_{(x,y,z)=(x,y,S(x,y))}$と勘違いしてしまう。
            \newline
            \textcolor{red}{$\partDerivLong{A_y(x,y,S(x,y))}{x}{}$と$\left.\partDeriv{A_y}{x}{}\right|_{(x,y,z)=(x,y,S(x,y))}$は別物}であることに注意せなばならぬ。
            連鎖律に注意して
            \begin{align*}
                \ctext{C} &= \left.\partDeriv{A_y}{x}{}\right|_{(x,y,z)=(x,y,S(x,y))}\partDeriv{x}{x}{} + \left.\partDeriv{A_y}{y}{}\right|_{(x,y,z)=(x,y,S(x,y))}\partDeriv{y}{x}{} + \left.\partDeriv{A_y}{z}{}\right|_{(x,y,z)=(x,y,S(x,y))}\partDeriv{S}{x}{} \\
                &= \partDeriv{A_y}{x}{} + \partDeriv{A_y}{z}{}\partDeriv{S}{x}{}
            \end{align*}
            同様にして
            \[\ctext{D} = \partDeriv{A_x}{y}{} + \partDeriv{A_x}{z}{}\partDeriv{S}{y}{}\]
            よって結局
            \begin{align*}
                \ctext{A} &= \iint_\Omega{\bracks*{\partDeriv{A_y}{x}{} + \partDeriv{A_y}{z}{}\partDeriv{S}{x}{} - \partDeriv{A_x}{y}{} - \partDeriv{A_x}{z}{}\partDeriv{S}{y}{}}\dx\dy} \\
                &= \iint_\Omega{\bracks*{\parens*{\partDeriv{A_y}{x}{}-\partDeriv{A_x}{y}{}} + \partDeriv{A_y}{z}{}\partDeriv{S}{x}{} - \partDeriv{A_x}{z}{}\partDeriv{S}{y}{}}\dx\dy}
            \end{align*}
            次に\ctext{B}を見る。$t$による偏微分を$\dot{}$で表すことにする。線積分の定義に従って
            \begin{align*}
                \ctext{B} &= \integrate{0}{1}{A_z(\phi(t),\psi(t),S(\phi(t),\psi(t))\parens*{\partDeriv{S}{x}{}\dot{\phi(t)}+\partDeriv{S}{y}{}\dot{\psi(t)}}}{}{t} \\
                &= \integrate{0}{1}{\parens*{A_z(\phi(t),\psi(t),S(\phi(t),\psi(t))\partDeriv{S}{x}{}}\dot{\phi(t)}}{}{t} + \integrate{0}{1}{\parens*{A_z(\phi(t),\psi(t),S(\phi(t),\psi(t))\partDeriv{S}{y}{}}\dot{\psi(t)}}{}{t} \\
                &= \oint_{\partial \Omega}{A_z(x,y,S(x,y))\partDeriv{S}{x}{}\dx} + \oint_{\partial \Omega}{A_z(x,y,S(x,y))\partDeriv{S}{y}{}\dy} \\
                &= \oint_{\partial \Omega}{\parens*{A_z(x,y,S(x,y))\partDeriv{S}{x}{}\dx + A_z(x,y,S(x,y))\partDeriv{S}{y}{}\dy}}
            \end{align*}
            これにGreenの定理を適用して，連鎖律に注意して偏微分していく。
            \begin{align*}
                \ctext{B} &= \iint_\Omega{\bracks*{\partDerivLong{\parens*{A_z(x,y,S(x,y))\partDeriv{S}{y}{}}}{x}{} - \partDerivLong{\parens*{A_z(x,y,S(x,y))\partDeriv{S}{x}{}}}{y}{}}\dx\dy} \\
                &= \iint_\Omega{\parens*{\partDeriv{A_z}{x}{}\partDeriv{S}{y}{} + \partDeriv{A_z}{z}{}\partDeriv{S}{x}{}\partDeriv{S}{y}{} + A_z\partDerivIIHetero{S}{y}{x} - \partDeriv{A_z}{y}{}\partDeriv{S}{x}{} - \partDeriv{nA_z}{z}{}\partDeriv{S}{y}{}\partDeriv{S}{x}{} - A_z\partDerivIIHetero{S}{x}{y}}\dx\dy} \\
                &= \iint_\Omega{\parens*{\partDeriv{A_z}{x}{}\partDeriv{S}{y}{} - \partDeriv{A_z}{y}{}\partDeriv{S}{x}{}}\dx\dy}\quad\parens*{\because\;\text{$S$は$C^2$級と仮定したから$\partDerivIIHetero{S}{y}{x} = \partDerivIIHetero{S}{x}{y}$}}
            \end{align*}
            以上より
            \begin{align*}
                \oint_c{A\cdot\dvc} &= \iint_\Omega{\bracks*{-\parens*{\partDeriv{A_z}{y}{}-\partDeriv{A_y}{z}{}}\partDeriv{S}{x}{}-\parens*{\partDeriv{A_x}{z}{}-\partDeriv{A_z}{x}{}}\partDeriv{S}{y}{}+\parens*{\partDeriv{A_y}{x}{}-\partDeriv{A_x}{y}{}}}\dx\dy}
            \end{align*}
            \par
            これが定理の右辺と等しいことを示せば良い。
            まず$\dS$の正体を暴こう。感覚的には，$\dS$とは，$S$上で$x,y$が微小量$\dx,\dy$だけ動いた時に舐める$S$の面積をノルムに持ち，向きはその微小面上の任意の代表点での上向きの法線ベクトルに一致するようなベクトルであった。
            所謂「面積ベクトル」である。
            \par
            厳密にはRiemann和を取る段階で定義され，$S$上の点$(x,y,S(x,y))$における接平面のうち，その射影が$xy$平面上の矩形$\dx\dy$に収まる部分(3次元上の平行四辺形)の面積をノルムに持ち，向きは接平面の上向き法線ベクトルと一致するようなベクトルである。
            \[\dS = \parens*{\vix\dx+\viz\partDeriv{S}{x}{}\dx}\times\parens*{\viy\dy+\viz\partDeriv{S}{y}{}\dy} = \parens*{-\vix\partDeriv{S}{x}{} - \viy\partDeriv{S}{y}{} + \viz}\dx\dy\]
            よって定理の右辺は
            \begin{align*}
                \iint_S{\nabla\times A}\cdot\dS &= \iint_\Omega{\bracks*{\vix\parens*{\partDeriv{A_z}{y}{}-\partDeriv{A_y}{z}{}}+\viy\parens*{\partDeriv{A_x}{z}{}-\partDeriv{A_z}{x}{}}+\viz\parens*{\partDeriv{A_y}{x}{}-\partDeriv{A_x}{y}{}}}\cdot\dS} \\
                &= \iint_\Omega{\bracks*{-\parens*{\partDeriv{A_z}{y}{}-\partDeriv{A_y}{z}{}}\partDeriv{S}{x}{}-\parens*{\partDeriv{A_x}{z}{}-\partDeriv{A_z}{x}{}}\partDeriv{S}{y}{}+\parens*{\partDeriv{A_y}{x}{}-\partDeriv{A_x}{y}{}}}\dx\dy} \\
                &= \oint_c{A\cdot\dvc}
            \end{align*}
        \end{proof}
    \section{諸注意}
        \subsection{ベクトル場の回転は元のベクトル場と直交するとは限らない}
            例：$\vA = (2r_1)\bm{i}_1 + \bm{i}_3$